\documentclass[a4paper,10pt,twoside,openany]{article}

\usepackage[lang=hebrew]{maths}
\usepackage{hebrewdoc}
\usepackage{stylish}
\usepackage{lipsum}
\let\bs\blacksquare

\usepackage{siunitx,array}
\sisetup{table-format=-1.5, group-digits=false}
\newcolumntype{L}{>{\phantom{$-$}}l}       % prefix some whitespace
\newcommand\mcL[1]{\multicolumn{1}{L}{#1}} % handy shortcut macro

\setlength{\parindent}{0pt}

%%%%%%%%%%%%
% Styling %
%%%%%%%%%%%%

\usepackage{enumitem}

%%%%%%%%%%%%%
% Counters  %
%%%%%%%%%%%%%

\setcounter{section}{1}     
            
%BIBLIOGRAPHY
\usepackage[
backend=biber,
style=alphabetic,
]{biblatex}
\addbibresource{bibliography.bib} %Imports bibliography file

%%%%%%%%%%
% Title  %
%%%%%%%%%%
\title{
אלגברה ב' - גיליון תרגילי בית 4 \\
הפולינום האופייני, ולכסינות
\\
\vspace{1cm}
\large{תאריך הגשה: 28.5.2026}
}
\date{}

\begin{document}
\maketitle

\begin{exercise}
עבור כל אחת מההעתקות הבאות $T \in \End\prs{V}$, קיבעו האם היא לכסינה. עבור אלו שהינן לכסינות, מיצאו בסיס
$\mcal{B}$
של $V$
ומטריצה אלכסונית
$D$
עבורם
$\brs{T}_\mcal{B} = D$. עבור אלו שאינן לכסינות, הוכיחו את טענתכן.

\begin{enumerate}
\item
\begin{align*}
T \colon \RR^2 &\to \RR^2 \\
\pmat{x \\ y} &\mapsto \pmat{-y \\ x}
\end{align*}

\item
\begin{align*}
T \colon \CC^2 \to \CC^2 \\
\pmat{x \\ y} &\mapsto \pmat{-y \\ x}
\end{align*}

\item
\begin{align*}
T \colon \Mat_2\prs{\RR} &\to \Mat_2\prs{\RR} \\
A &\mapsto A^t
\end{align*}

\item
\begin{align*}
T \colon \CC_2\brs{x} &\to \CC_2\brs{x} \\
p\prs{x} &\mapsto p'\prs{x}
\end{align*}
\end{enumerate}
\end{exercise}

\begin{exercise}
תהי
$A \in \Mat_n\prs{\FF}$
עם פולינום אופייני
\begin{align*}
\text{.} p_A\prs{x} = \sum_{i=0}^n a_i x^i
\end{align*}
הוכיחו כי
$a_{n-1} = -\tr\prs{A}$.
\end{exercise}

\begin{exercise}
יהי
$T \in \End_{\FF}\prs{V}$
אופרטור לינארי, ויהי
$\lambda \in \FF$.
הוכיחו כי
\begin{align*}
\text{.} p_{\lambda T}\prs{x} = \lambda^{\dim_\FF\prs{V}} p_T\prs{\frac{x}{\lambda}}
\end{align*}
\end{exercise}

\begin{exercise}
תהיינה
$A,B \in \Mat_n\prs{\FF}$.
הוכיחו או הפריכו את הטענות הבאות. עבור הטענות שאינן נכונות, מיצאו דוגמה נגדית.

\begin{enumerate}
\item $p_{AB}\prs{x} = p_{A}\prs{x} p_{B}\prs{x}$.
\item $p_{AB}\prs{x} = p_{BA}\prs{x}$.
\item $p_{A+B}\prs{x} = p_A\prs{x} + p_B\prs{x}$.
\item אם
$A,B$
לכסינות, גם
$AB$
לכסינה.
\item אם
$A,B$
לכסינות, גם
$A+B$
לכסינה.
\end{enumerate}
\end{exercise}

\end{document}